% theme: inscribed_angle_theorem
% 表：円周角・中心角、直径に対する円周角。裏：弧との関係、円に内接する四角形。
\MajorQuestionLesson{円周角の定理}
\SetRowSpace{2mm}

\TypeQuestion{点 $\mathrm{O}$ は円の中心である。図の $x$ の大きさを求めなさい。}{inscribed_angle_and_central_angle}
\QQFigRowThree[30mm]{fig/lesson_01_central_01}{}{fig/lesson_01_central_02}{}{fig/lesson_01_central_03}{}
\QQFigRowThree[30mm]{fig/lesson_01_central_04}{}{fig/lesson_01_central_05}{}{fig/lesson_01_central_06}{}

\TypeQuestion{太線で示した線分は円の直径である。図の $x$ の大きさを求めなさい。}{diameter_and_right_angle}
\QQFigRowThree[30mm]{fig/lesson_01_diameter_01}{}{fig/lesson_01_diameter_02}{}{fig/lesson_01_diameter_03}{}
\QQFigRowThree[30mm]{fig/lesson_01_diameter_04}{}{fig/lesson_01_diameter_05}{}{fig/lesson_01_diameter_06}{}

\AnswerColumnBreak
\LessonQuestionBreak

\TypeQuestion{図の数値と印を用いて、$x$ の大きさを求めなさい。同じ印をつけた弧は等しい。}{arc_and_inscribed_angle}
\noindent
\QQFigTextRowThreeCell[25mm]{fig/lesson_01_arc_01}{弧 $\mathrm{AB}:$ 弧 $\mathrm{BC}=2:3$、$\Ang{ADB}=40^\circ$、$x=\Ang{BEC}$}{}
\hfill
\QQFigTextRowThreeCell[25mm]{fig/lesson_01_arc_02}{弧 $\mathrm{AB}=$ 弧 $\mathrm{CD}$、$\Ang{ADB}=32^\circ$、$x=\Ang{CED}$}{}
\hfill
\QQFigTextRowThreeCell[25mm]{fig/lesson_01_arc_03}{弧 $\mathrm{AB}:$ 弧 $\mathrm{CD}=3:5$、$\Ang{ADB}=36^\circ$、$x=\Ang{CED}$}{}
\QQFigTextRowThreeEnd
\noindent
\QQFigTextRowThreeCell[25mm]{fig/lesson_01_arc_04}{円周は $24\mathrm{cm}$、弧 $\mathrm{AB}$ は $8\mathrm{cm}$、$x=\Ang{ADB}$}{}
\hfill
\QQFigTextRowThreeCell[25mm]{fig/lesson_01_arc_05}{弧 $\mathrm{AB}=$ 弧 $\mathrm{BC}$、$\Ang{ADB}=54^\circ$、$x=\Ang{BEC}$}{}
\hfill
\QQFigTextRowThreeCell[25mm]{fig/lesson_01_arc_06}{弧 $\mathrm{AB}:$ 弧 $\mathrm{CD}=2:5$、$\Ang{ADB}=36^\circ$、$x=\Ang{CED}$}{}
\QQFigTextRowThreeEnd

\TypeQuestion{四角形 $\mathrm{ABCD}$ は円に内接している。図の $x$ の大きさを求めなさい。}{cyclic_quadrilateral_angles}
\QQFigRowThree[30mm]{fig/lesson_01_cyclic_01}{}{fig/lesson_01_cyclic_02}{}{fig/lesson_01_cyclic_03}{}
\QQFigRowThree[30mm]{fig/lesson_01_cyclic_04}{}{fig/lesson_01_cyclic_05}{}{fig/lesson_01_cyclic_06}{}
